\documentclass[a4paper,11pt]{article}
\usepackage[T1]{fontenc}
\usepackage[utf8]{inputenc}
\usepackage{amsmath}
\usepackage{siunitx}
\usepackage[margin=2.5cm]{geometry}

\title{ATCG I Sheet R01 Solution}
\author{Radiometry and Tonemapping}
\date{April 23, 2026}

\begin{document}
\maketitle

\section*{Assignment 3: Radiance and Irradiance}
The table is illuminated by the sun, modeled as a distant emitter with constant radiance
\[
L_s = 20.045\,\si{\mega\watt\per\square\meter\per\steradian}
= 2.0045 \times 10^7\,\si{\watt\per\square\meter\per\steradian}.
\]
The table normal forms an angle of $\theta = 45^\circ$ with the incident direction, and the table size is
\[
A = 0.8 \times 0.8 = 0.64\,\si{\square\meter}.
\]

The irradiance from a source of constant radiance is
\[
E = \int_{\Omega_s} L_s \cos\theta \, d\omega
= L_s \cos\theta \int_{\Omega_s} d\omega
= L_s \cos\theta \, \Omega_s.
\]
So we first need the solid angle subtended by the sun. Using the solar radius
\[
R_s = 6.9634 \times 10^8\,\si{\meter}
\]
and the mean Earth-sun distance
\[
D = 1.495978707 \times 10^{11}\,\si{\meter},
\]
the angular radius is
\[
\alpha = \arcsin\left(\frac{R_s}{D}\right).
\]
The exact solid angle of a disk with angular radius $\alpha$ is
\[
\Omega_s = 2\pi(1-\cos\alpha).
\]
Numerically,
\[
\alpha \approx \arcsin(4.6547 \times 10^{-3}),
\qquad
\Omega_s \approx 6.8068 \times 10^{-5}\,\si{\steradian}.
\]

Therefore the irradiance at the table is
\begin{align*}
E
&= L_s \cos(45^\circ)\,\Omega_s \\
&= 2.0045 \times 10^7 \cdot \frac{1}{\sqrt{2}} \cdot 6.8068 \times 10^{-5} \\
&\approx 9.648 \times 10^2\,\si{\watt\per\square\meter}.
\end{align*}
Thus,
\[
\boxed{E \approx 964.8\,\si{\watt\per\square\meter}}.
\]

The total radiant power incident on the table is
\[
\Phi = E \cdot A = 964.8 \times 0.64 \approx 617.5\,\si{\watt}.
\]
Hence,
\[
\boxed{\Phi \approx 617.5\,\si{\watt}}.
\]

\section*{Practical Assignments}
Assignments 1 and 2 are solved in the framework code:
\begin{itemize}
\item Robertson HDR reconstruction is implemented in:
\[
\texttt{src/exercise01\_HDR/Robertson.cpp}
\quad \text{and} \quad
\texttt{src/exercise01\_HDR/Robertson.cu}
\]
\item Tone mapping modes \texttt{linear\_max}, \texttt{linear\_fixed}, \texttt{gamma\_fixed}, and \texttt{histogram} are implemented in:
\[
\texttt{src/exercise01\_HDR/Tonemapping.cpp}
\quad \text{and} \quad
\texttt{src/exercise01\_HDR/Tonemapping.cu}
\]
\end{itemize}

\end{document}
